% Standalone section fragment for the gradient-dependent curvature model.
% No optimality claim is made for the new explicit certificate.
\section{Boundary minima and a computable error bound}
\label{sec:boundary-gradient-model}

Let $Q$ be a nontrivial compact convex set, let $f$ be $C^2$ near $Q$,
and let $x^*\in\arg\min_Q f$. Assume, for every $z\in Q$, that
\begin{equation}\label{eq:boundary-gradient-assumptions}
\begin{aligned}
 \langle\nabla f(z)-\nabla f(x^*),z-x^*\rangle
    &\ge m\|z-x^*\|^2,\qquad m>0,\\
 \|\nabla^2 f(z)\|&\le L_0+L_1\|\nabla f(z)\|,
       \qquad L_0,L_1\ge0.
\end{aligned}
\end{equation}
Write $D\ge\operatorname{diam}Q>0$ and
\[
 G(x)=\max_{v\in Q}\langle\nabla f(x),x-v\rangle.
\]
The first condition is shifted restricted secant regularity. A boundary
minimizer may have a nonzero gradient. The second condition controls
curvature by the gradient norm; it does not require this norm to vanish.

\subsection{What the available observations can certify}

\begin{proposition}[A limitation of the gradient norm]
\label{prop:boundary-norm-obstruction}
Fix $m,D>0$. Suppose a certificate $\mathcal C$ uses only the gradient
norm and the constants $m,L_0=m,L_1=0,D$, and is valid for every problem
satisfying \eqref{eq:boundary-gradient-assumptions}:
\[
 f(x)-f(x^*)\le\mathcal C(\|\nabla f(x)\|).
\]
Then, for every scalar argument $t\ge0$,
\begin{equation}\label{eq:boundary-norm-lower}
 \mathcal C(t)\ge
 \begin{cases}
 t^2/(2m),&0\le t\le mD,\\
 tD-mD^2/2,&t\ge mD.
 \end{cases}
\end{equation}
Consequently, such a certificate can stay arbitrarily large at an exact
minimizer, even with constant positive curvature.
\end{proposition}
\begin{proof}
On the fixed interval $[0,D]$, consider
$f_a(x)=ax+mx^2/2$, where $a\ge0$. Its minimum is at zero, its Hessian
is $m$, and all the stated constants are independent of $a$.
At a point with derivative $t$, one has $a=t-mx$ and
\[
 0\le x\le\min\{D,t/m\},\qquad
 f_a(x)-f_a(0)=tx-mx^2/2.
\]
This expression increases throughout the indicated interval. Its
maximum is exactly the right side of \eqref{eq:boundary-norm-lower}.
The two cases are attained by $a=0,x=t/m$ and $a=t-mD,x=D$,
respectively. At the exact minimizer of $f_a$, the norm is $a$ and the
error is zero, but the same valid function $\mathcal C(a)$ must satisfy
\eqref{eq:boundary-norm-lower}.
\end{proof}

\begin{proposition}[The full gap alone has no uniform error bound]
\label{prop:boundary-gap-obstruction}
Under \eqref{eq:boundary-gradient-assumptions} with $L_1>0$, the
values of $G(x),m,L_0,L_1,D$ do not determine a finite universal upper
bound on $f(x)-f(x^*)$. This remains true on the fixed square $[0,1]^2$,
even when the curvature condition holds on all of $\mathbb R^2$.
\end{proposition}
\begin{proof}
Take $Q=[0,1]^2$, $m=1/2$, and $L_0=L_1=1$. For $t\ge0$, define
\[
 f_t(u,v)=\frac{u^2}{2}+\frac t\pi(1-\cos(\pi u))
             +\sqrt{1+v^2}+(\pi t+1)v.
\]
Its gradient is
$(u+t\sin(\pi u),\,\pi t+1+v/\sqrt{1+v^2})$.
Both coordinates are nonnegative on $Q$, so $x^*=(0,0)$ is a minimizer.
Moreover, on $Q$,
\[
 \langle\nabla f_t(u,v)-\nabla f_t(0,0),(u,v)\rangle
 =u^2+tu\sin(\pi u)+\frac{v^2}{\sqrt{1+v^2}}
 \ge\tfrac12(u^2+v^2).
\]
On all of $\mathbb R^2$, the Hessian is
$\operatorname{diag}(1+\pi t\cos(\pi u),(1+v^2)^{-3/2})$,
and the second gradient coordinate is at least $\pi t$. Hence
\[
 \|\nabla^2 f_t(u,v)\|\le1+\pi t
 \le1+\|\nabla f_t(u,v)\|.
\]
All constants and the domain are fixed. At $x=(1,0)$, the linear
minimization oracle selects $(0,0)$, and
\[
 G(1,0)=1,\qquad
 f_t(1,0)-f_t(0,0)=\frac12+\frac{2t}{\pi}.
\]
The latter tends to infinity with $t$ while the former stays fixed.
\end{proof}

This example allows negative curvature away from the minimizer, while
retaining shifted secant regularity. A valid error bound therefore needs
information beyond the full gap and the stated constants. The current
gradient norm provides one such quantity.

\subsection{A local curvature bound and an explicit certificate}

\begin{theorem}[A certificate without a vanishing gradient norm]
\label{thm:boundary-local-certificate}
Under \eqref{eq:boundary-gradient-assumptions}, set
\begin{align}
 R(x)&=\min\left\{D,\sqrt{G(x)/m},\|\nabla f(x)\|/m\right\},
       \label{eq:boundary-local-radius}\\
 K(x)&=\max\left\{m,
    \bigl(L_0+L_1\|\nabla f(x)\|\bigr)e^{L_1R(x)}\right\},
       \label{eq:boundary-local-curvature}\\
 C(x)&=G(x)+
   \max\left\{0,\frac{K(x)^2-2mK(x)-m^2}{4K(x)}\right\}R(x)^2.
       \label{eq:boundary-local-certificate}
\end{align}
Then $f(x)-f(x^*)\le C(x)$ for every $x\in Q$.
For every sequence in this fixed problem with $G(x_k)\to0$, one has
$C(x_k)\to0$. No condition $\|\nabla f(x_k)\|\to0$ is needed.
\end{theorem}
\begin{proof}
For $x\ne x^*$, let $r=\|x-x^*\|$ and $e=(x-x^*)/r$.
Constrained optimality gives
$\langle\nabla f(x^*),e\rangle\ge0$. Shifted secant regularity gives
\[
 \langle\nabla f(x),e\rangle\ge mr,\qquad
 G(x)\ge r\langle\nabla f(x),e\rangle\ge mr^2.
\]
Cauchy--Schwarz also gives $mr\le\|\nabla f(x)\|$, hence $r\le R(x)$.

Along the segment followed backwards from $x$, the norm is absolutely
continuous and satisfies, almost everywhere,
\[
 \frac{d}{ds}\|\nabla f(x-se)\|
 \le L_0+L_1\|\nabla f(x-se)\|.
\]
The scalar differential inequality implies
\[
 L_0+L_1\|\nabla f(x-se)\|
 \le\bigl(L_0+L_1\|\nabla f(x)\|\bigr)e^{L_1s},
       \qquad 0\le s\le r.
\]
For $L_1=0$ the same conclusion is immediate. Thus $K(x)$ bounds
the Hessian norm on $[x^*,x]$.

Define $u(s)=f(x^*+se)-f(x^*)$ on $[0,r]$, and abbreviate $K=K(x)$.
The preceding facts give $|u''|\le K$ and
$u'(0)\le u'(r)-mr$. Consequently
\[
 u'(s)\le
 \min\{u'(r)-mr+Ks,\ u'(r)+K(r-s)\}.
\]
Since $K\ge m$, the two affine bounds meet at
$s=(K+m)r/(2K)\in[0,r]$. Integrating the two pieces gives
\begin{equation}\label{eq:boundary-tent-integral}
 u(r)\le r u'(r)
       +\frac{K^2-2mK-m^2}{4K}r^2.
\end{equation}
Now $r u'(r)\le G(x)$. If the coefficient of $r^2$ is nonpositive,
discard that term. Otherwise use $r^2\le R(x)^2$.
This proves \eqref{eq:boundary-local-certificate}. If $x=x^*$,
constrained optimality gives $G(x)=0$, and the claim holds directly.

Finally, continuity of the gradient on compact $Q$ bounds its norm.
Hence $K(x_k)$ remains bounded, while $G(x_k)\to0$ implies
$R(x_k)\to0$. Formula \eqref{eq:boundary-local-certificate} then tends
to zero.
\end{proof}

\paragraph{Interpretation and limits.}
The rule $C(x)\le\varepsilon$ certifies the requested error.
The theorem proves convergence of the certificate along sequences with
vanishing full gap; convergence of a particular algorithm requires its
own argument. Large gradients can make this bound conservative. If a
verified bound $\|\nabla^2f(z)\|\le K_Q$ on $Q$ is also available,
$K(x)$ may be replaced by $\min\{K(x),K_Q\}$. The counterexample above
precludes a uniform error calibration based only on
$m,L_0,L_1,D$ and the gap. Quantitative work estimates obtained from
such a calibration need additional information, such as a gradient
or curvature bound and an initial error bound. The example is not an
iteration lower bound. The explicit certificate is not asserted to be
sharp for the gradient-dependent curvature class.
